Enterprise systems face a structural problem: the surfaces through which change
enters the system are too many, too permeable, and too decoupled from the
processes that are supposed to govern them. Consider three incidents, each drawn
from a pattern that recurs across industries:

\begin{description}[leftmargin=1.5em]
  \item[Schema-valid fraud.] A payment-operations employee with database access
        issues \texttt{UPDATE invoice SET status = 'paid' WHERE id = 4711}. The
        row passes every schema constraint --- the column accepts the value, the
        foreign keys resolve, the audit trigger fires and dutifully records the
        write. Yet no approval event ever occurred. The fraud is invisible to
        schema validation because the \emph{shape} of the data is correct; only
        its \emph{process position} is wrong.
  \item[Rogue administrator.] A platform administrator, asked to ``unblock'' a
        stalled onboarding case, sets a user's role to \texttt{admin} directly in
        the identity store. The four-eyes approval workflow that exists precisely
        for this operation is bypassed, not by exploiting a vulnerability, but by
        using a privilege the architecture grants: the power to mutate state
        without emitting the events the process model requires.
  \item[AI-agent mutation.] An autonomous agent equipped with a write-capable
        API tool misinterprets an instruction and cancels forty active
        subscriptions. Every API call is authenticated and authorized; the agent
        holds a valid service token. What is missing is any binding between the
        agent's actions and a process law that says \emph{from this state, with
        this mandate, that transition is admissible}.
\end{description}

In all three cases the failure is the same: \textbf{the system accepted a change
because it could be parsed and authorized, not because it could be proven to be
a lawful step of a declared process.} Current defences respond after the fact:
logs are written, anomaly detectors flag deviations~\cite{ko_comuzzi_anomaly_2021},
and audit teams attempt to reconstruct what happened, often months later.
Process mining, in van der Aalst's canonical formulation, discovers and verifies
process models against event logs recorded after
execution~\cite{van_der_aalst_2016_process_mining}. The discipline is analytic
and retrospective by construction.

The thesis of this paper is that the temporal position of that discipline should
be inverted. Instead of mining a log to determine whether a completed process
conformed, we require that a \emph{process-valid, cryptographically witnessed
event} be produced and admitted \emph{before} the change is permitted to alter
enterprise state. Under this regime the unit of enterprise change is no longer
the database write or the API call; it is the \emph{admitted change}: a tuple
binding actor, event type, object relations, current process state, prior
receipt, policy version, validator version, freshness nonce, and a post-quantum
signature over all of them. Anything that cannot produce this witness is not
interpreted as valid change. It is refused, or it is \emph{residualized} ---
quarantined as evidence without state authority.

We do not claim such a system is unhackable. We claim something narrower and,
we argue, more useful: that under explicitly named cryptographic and
implementation assumptions, entire \emph{classes} of adversarial mutation become
\emph{mathematically inaccessible to acceptance}. The adversary may still submit
anything; what they cannot do --- without breaking a named assumption --- is make
it count.

\subsection{Contributions}

This paper makes four contributions:

\begin{enumerate}[label=(\arabic*),leftmargin=2em]
  \item \textbf{The Enterprise Change Admissibility Framework}
        (Section~\ref{sec:model}). We define the enterprise change universe
        $\Omega_{\mathrm{change}}$ and the admitted change set $\Aent$ as the
        subset satisfying a seven-conjunct acceptance predicate. We give the
        state transition law under which non-admitted changes are identity
        transitions, prove the membership theorem, and formalise
        \emph{residualization} as the lossless refusal mode. The framework
        extends boundary predicates and admitted symbolic events from process
        mining theory~\cite{van_der_aalst_2016_process_mining,van_der_aalst_berti_2020}
        into a governance law for all enterprise state change.
  \item \textbf{Eight inaccessibility propositions}
        (Section~\ref{sec:inaccessibility}). We prove that raw database writes,
        process-state skips, object substitution, backdating, replay,
        schema-valid fraud, rogue administrator edits, and AI-agent hallucinated
        actions are each inaccessible to acceptance, identifying for every class
        the specific predicate conjunct that blocks it. A running
        procure-to-pay example grounds each proposition in a concrete enterprise
        scenario.
  \item \textbf{A named attack-surface decomposition}
        (Section~\ref{sec:attacks}). We bound the probability of accepting an
        invalid change by a union over seven named failure channels ---
        cryptographic, algorithmic, modelling, operational, and
        implementation --- each with a characterised mitigation. This replaces
        unfalsifiable security claims (``zero trust,'' ``secure by design'')
        with an auditable assumption ledger, and connects each channel to the
        maturity-model literature for
        operationalisation~\cite{sok_ccmm_2024,crypto_agility_maturity_2022}.
  \item \textbf{A reference implementation and evaluation}
        (Section~\ref{sec:impl}). We instantiate the framework in
        \emph{wasm4pm}~\cite{chatman_wasm4pm_2026}, a process mining platform
        whose 60 algorithms compile from Rust to deterministic WebAssembly.
        Admitted changes are realised as \emph{TrueX envelopes} over canonical
        OCEL~2.0 logs~\cite{berti_et_al_ocel2_spec_2024}; receipts chain under
        BLAKE3; the boundary predicate is enforced as thirteen non-overlapping
        refusal states; and signing is post-quantum
        (ML-DSA~\cite{nist_fips_204_2024}). We report on the proof-gate test
        suite, determinism guarantees, and the operational verification
        protocol.
\end{enumerate}

\subsection{Running Example}
\label{sec:running-example}

Throughout the paper we use a procure-to-pay (P2P) process as the running
example. An \textsf{Invoice} object progresses through states
\[
  \textsf{draft} \;\to\; \textsf{submitted} \;\to\; \textsf{approved}
  \;\to\; \textsf{paid} \;\to\; \textsf{archived},
\]
driven by event types \textsf{Submit}, \textsf{Approve}, \textsf{Pay}, and
\textsf{Archive}. Each event relates the invoice to other objects: the
\textsf{Employee} who acts, the \textsf{PurchaseOrder} it references, and the
\textsf{BankTransaction} that settles it. The declared process law requires
that \textsf{Approve} be performed by an actor distinct from the submitter
(four-eyes), that \textsf{Pay} reference an approved invoice and a matching
purchase order (three-way match), and that no state be skipped. P2P is chosen
because it is the canonical target of enterprise fraud: every attack class in
Section~\ref{sec:inaccessibility} has a documented real-world instance in this
process.

\subsection{Structure}

Section~\ref{sec:background} reviews process mining, OCEL~2.0, receipt chains,
and post-quantum signing. Section~\ref{sec:model} develops the formal model,
including the threat model, and proves the membership theorem.
Section~\ref{sec:inaccessibility} proves the eight inaccessibility propositions.
Section~\ref{sec:attacks} gives the attack-surface decomposition and discusses
residual risk. Section~\ref{sec:impl} describes and evaluates the wasm4pm
implementation. Section~\ref{sec:related} positions the work against eight
adjacent research threads. Section~\ref{sec:conclusion} concludes with
limitations and the category claim.
